栈、单调栈和队列深度题卡 按题型拆成章内大节，但每个大节都先从 LeetCode 案例切入，再进入分流、案例矩阵、案例推导和实验复盘。这样读者看到的是从具体题推到规律，而不是先读抽象方法再猜它能用在哪。

\topic{栈、单调栈与表达式}
这一节先看 LeetCode 20 有效的括号、LeetCode 32 最长有效括号、LeetCode 42 接雨水 这几道 LeetCode 题，再整理同一题型的分流和推导。阅读顺序不是先背原理，而是先看案例的暴力瓶颈，再把重复出现的观察抽象成方法。

\subsection{原理和适用条件}

以 LeetCode 20 有效的括号、LeetCode 32 最长有效括号、LeetCode 42 接雨水 为主线：LeetCode 20 有效的括号：模型是“括号匹配要求最近打开的左括号先被关闭。”，优化抓手是“栈保存尚未匹配的左括号，遇到右括号必须和栈顶类型一致。”；LeetCode 32 最长有效括号：模型是“有效括号子串需要知道每段非法边界之后的最早可连接位置。”，优化抓手是“栈保存下标，先放入哨兵负一；匹配后用当前下标减栈顶得到长度。”；LeetCode 42 接雨水：模型是“每个位置的水由左侧最高和右侧最高的较小者决定。”，优化抓手是“前后缀、双指针、单调栈都是不同状态维护方式；要能说明各自结算对象。”。这些案例共同推出的规律是：先对每个位置向左或向右扫描，再识别还没有被解决的候选。栈保存等待未来元素解决的对象，新元素到来时一次性结算被它支配的一批旧候选。 方法名只是复盘后的标签，真正要掌握的是从题面约束、暴力失败、状态设计到边界反例的推导链。

\begin{principlebox}{图示：从 LeetCode 案例到 栈、单调栈与表达式推导链}
\begin{tabular}{@{}p{0.18\linewidth}p{0.74\linewidth}@{}}
暴力基线 & 枚举候选、完整模拟或逐个检查，先保证覆盖全部可能答案。\\
瓶颈定位 & 瓶颈是未解决元素没有被组织起来。单调栈保存的是仍在等待某个未来元素解决的候选；普通栈保存的是最近打开但尚未关闭的结构。\\
结构选择 & 当新元素到来时，它会一次性解决栈顶一串被它支配的旧元素。被弹出的元素以后再也不会参与比较，因此总体线性。\\
不变量 & 栈内下标对应的值要保持单调，或者栈顶必须是最近未匹配对象。弹出时要说清楚当前元素充当右边界、下一个栈顶充当左边界，还是当前运算符触发计算。\\
代码落地 & 栈里通常存下标而不是值，因为要计算距离、宽度和过期。可以用 \code{std::vector<int>} 当栈，便于查看顶部和减少适配器限制。\\
\end{tabular}
\end{principlebox}

\subsection{LeetCode 案例分流}

\begin{principlebox}{从 LeetCode 案例分流：栈、单调栈与表达式}
\textbf{单调栈求下一个边界。}
代表案例：LeetCode 739 每日温度、LeetCode 496 下一个更大元素 I、LeetCode 503 下一个更大元素 II。先看这些题的题面条件：每个位置要找左侧或右侧第一个更大、更小元素。由此推出的处理方式是：栈保存尚未解决的下标，新元素到来时弹出并结算。风险点：相等元素如何处理会影响答案和去重。

\textbf{宽度和面积结算。}
代表案例：LeetCode 84 柱状图最大矩形、LeetCode 85 最大矩形、LeetCode 42 接雨水。先看这些题的题面条件：某个元素作为最矮或最低边界时，需要左右第一个破坏条件的位置。由此推出的处理方式是：弹出时用当前下标作右边界、弹出后栈顶作左边界。风险点：哨兵和宽度公式是主要错误来源。

\textbf{括号和表达式上下文。}
代表案例：LeetCode 20 有效括号、LeetCode 32 最长有效括号、LeetCode 224 基本计算器、LeetCode 227 基本计算器 II、LeetCode 394 字符串解码。先看这些题的题面条件：题目存在嵌套结构、最近未闭合结构或运算优先级。由此推出的处理方式是：栈保存上一层状态、未匹配括号或已经归约的项。风险点：多位数、空栈、操作数顺序和括号结束时机要单独测试。

\textbf{双栈和频率栈。}
代表案例：LeetCode 232 用栈实现队列、LeetCode 895 最大频率栈。先看这些题的题面条件：需要用栈模拟另一种顺序结构，或同时按频率和最近性弹出。由此推出的处理方式是：用多个栈分别保存输入、输出或同频元素。风险点：均摊复杂度和多索引一致性要讲清楚。

\end{principlebox}

\subsection{案例矩阵}

本节案例覆盖 LeetCode 20 有效的括号、LeetCode 32 最长有效括号、LeetCode 42 接雨水、LeetCode 71 简化路径、LeetCode 84 柱状图中最大的矩形、LeetCode 85 最大矩形、LeetCode 150 逆波兰表达式求值、LeetCode 155 最小栈、LeetCode 224 基本计算器、LeetCode 227 基本计算器 II、LeetCode 232 用栈实现队列、LeetCode 239 滑动窗口最大值、LeetCode 394 字符串解码、LeetCode 402 移掉 K 位数字、LeetCode 496 下一个更大元素 I、LeetCode 503 下一个更大元素 II、LeetCode 739 每日温度、LeetCode 895 最大频率栈。阅读时不要按题号背答案，而要先判断题面落在哪个分流场景：查询目标是什么，输入约束给了什么单调性或等价关系，暴力慢在哪里，优化结构保存了什么状态。

\begin{principlebox}{案例矩阵}
\textbf{LeetCode 20 有效的括号。}
括号匹配要求最近打开的左括号先被关闭。单点风险：空串、以右括号开头、交叉匹配、结束后栈非空。

\textbf{LeetCode 32 最长有效括号。}
有效括号子串需要知道每段非法边界之后的最早可连接位置。单点风险：空串、全左括号、全右括号、多个断开的合法段。

\textbf{LeetCode 42 接雨水。}
每个位置的水由左侧最高和右侧最高的较小者决定。单点风险：单调递增、单调递减、平台高度、多个凹槽、数组长度不足三。

\textbf{LeetCode 71 简化路径。}
Unix 绝对路径由斜杠分隔目录段，点号段和返回上级目录段需要被规范化。单点风险：连续斜杠、根目录上继续 ..、路径末尾斜杠、目录名含点但不是点号。

\textbf{LeetCode 84 柱状图中最大的矩形。}
每根柱子作为最矮高度时，需要知道左右第一个更矮位置。单点风险：递增、递减、相等高度、哨兵零、宽度计算。

\textbf{LeetCode 85 最大矩形。}
每一行都可以把连续的一当成柱状图高度。单点风险：全零、全一、单行、单列、宽度和高度结算。

\textbf{LeetCode 150 逆波兰表达式求值。}
后缀表达式把优先级和括号都编码进 token 顺序。单点风险：负数 token、多位数、除法向零截断、减法顺序。

\textbf{LeetCode 155 最小栈。}
普通栈只能看到栈顶，不能常数时间知道历史最小值。单点风险：重复最小值、弹出最小值、空栈操作约定。

\textbf{LeetCode 224 基本计算器。}
括号会开启新的表达式上下文，括号结束时要回到上一层。单点风险：多位数、前导空格、一元负号、嵌套括号。

\textbf{LeetCode 227 基本计算器 II。}
乘除优先级高于加减，可以在读到下一个运算符时结算上一项。单点风险：多位数、空格、除法向零截断、最后一个数字需要结算。

\textbf{LeetCode 232 用栈实现队列。}
两个栈分别负责输入顺序和输出顺序。单点风险：连续 peek、空队列、交替 push pop。

\textbf{LeetCode 239 滑动窗口最大值。}
每个窗口最大值来自仍未过期且没有被新元素支配的候选。单点风险：k 为一、k 等于数组长度、重复最大值、队首过期。

\textbf{LeetCode 394 字符串解码。}
嵌套结构需要保存上一层字符串和重复次数。单点风险：多位数字、嵌套、空片段、数字后必须跟括号。

\textbf{LeetCode 402 移掉 K 位数字。}
要删除 k 位使剩余数字最小，本质是让高位尽量小。单点风险：前导零、k 未用完、k 等于长度、数字已经递增。

\textbf{LeetCode 496 下一个更大元素 I。}
第二个数组提供每个值右侧第一个更大值的全局关系。单点风险：不存在更大值、递减数组、值唯一假设。

\textbf{LeetCode 503 下一个更大元素 II。}
环形数组可以通过遍历两倍下标模拟绕回。单点风险：全递减、全相等、长度一、取模下标。

\textbf{LeetCode 739 每日温度。}
每一天等待右侧第一个更高温度。单点风险：没有更高温、相等温度、递增、递减。

\textbf{LeetCode 895 最大频率栈。}
弹出时既要频率最高，又要在同频中最近入栈。单点风险：重复 push、pop 后最大频率下降、同频最近性、空栈不在题意内。

\end{principlebox}

\subsection{共性落点}

\begin{principlebox}{本组共性落点}
\textbf{慢版本。} 暴力做法对每个位置向左或向右寻找第一个满足条件的元素，或者反复删除已经匹配的局部结构。这会让同一个元素被比较很多次。
\textbf{瓶颈。} 瓶颈是未解决元素没有被组织起来。单调栈保存的是仍在等待某个未来元素解决的候选；普通栈保存的是最近打开但尚未关闭的结构。
\textbf{状态字段。} 栈内对象的业务含义、当前输入、弹出时结算的对象、左边界和右边界；栈中存下标还是值要明确。
\textbf{不变量。} 栈内下标对应的值要保持单调，或者栈顶必须是最近未匹配对象。弹出时要说清楚当前元素充当右边界、下一个栈顶充当左边界，还是当前运算符触发计算。
\textbf{实现检查。} 弹栈前确认栈非空，弹出后再读取新的左边界；相等元素和哨兵高度要有统一策略。
\textbf{C++ 落地。} 栈里通常存下标而不是值，因为要计算距离、宽度和过期。可以用 \code{std::vector<int>} 当栈，便于查看顶部和减少适配器限制。
\textbf{证明抓手。} 证明从弹出时机开始：被弹出的对象在当前时刻答案已经确定，未来元素无法再改变它。
\end{principlebox}

\subsection{案例推导}

\noindent\textbf{LeetCode 20 有效的括号。}

\textbf{题目说明。} LeetCode 20 有效的括号 的题面入口要先还原成具体任务：括号匹配要求最近打开的左括号先被关闭。

\textbf{样例入口。} 拿 ([)] 手算：只计数左右括号会误判，因为闭合顺序错了。栈顶必须保存最近的未匹配左括号；遇到 ] 时栈顶若不是 [，立即失败。

\textbf{暴力基线。} 暴力反复删除字符串中的 ()、[]、\{\}，直到不能再删除；若最后为空则合法。

\textbf{暴力过程。} 以 ([)] 为例，左右括号数量都对，但没有任何相邻合法括号对可删除，说明只数数量会误判顺序。

\textbf{重复工作。} 题面模型：括号匹配要求最近打开的左括号先被关闭。暴力版的慢点要落到本题：浪费在反复扫描字符串找可删除对。括号匹配只关心最近一个尚未关闭的左括号，它天然符合后进先出。后面的优化只消掉这类重复工作，不能改变答案集合。

\textbf{优化条件。} 本题的关键观察是：栈保存尚未匹配的左括号，遇到右括号必须和栈顶类型一致。这句话必须能翻译成可维护状态；如果题目条件改变后观察不再成立，就不能继续套原来的写法。

\textbf{相邻状态。} 规律是括号匹配不是数量问题，而是最近未闭合结构的后进先出问题。把这个特例继续拆开看：栈里的对象都是答案尚未确定的候选；一旦当前元素触发弹栈，被弹出的候选必须已经同时拿到了左边界和右边界，未来元素不再有资格改变它的答案。

\textbf{状态复用。} 栈保存未匹配的左括号。遇到左括号入栈；遇到右括号时栈不能为空，且栈顶必须是对应类型，匹配后弹出。手算时只改被新输入影响的那一小部分，而不是回头重算完整候选。

\textbf{指针和字段。} 状态字段要能说清：stack.top 是最近未匹配左括号，ch 是当前字符，pairs 记录右括号需要的左括号类型。手算时按这个协议跟踪：处理 ([)]：入栈 (，入栈 [；遇到 ) 时栈顶是 [，不是 (，立即失败。

\textbf{代码落点。} 关键代码是遇到右括号先检查 empty，再比较 top。扫描结束后还必须检查栈为空，防止 (( 这类前缀合法但未闭合的输入。关键代码行必须能逐句翻译成“旧状态怎样变成新状态”，否则就只是模板。

\textbf{正确性。} 合法括号串的任意右括号都必须关闭最近的未闭合左括号；栈顶正好表示这个对象，所以每次局部匹配都是必要条件。

\textbf{边界实验。} 先用空串、以右括号开头、交叉匹配、结束后栈非空做反例检查。实验时准备空串、以右括号开头、交叉匹配、结束后栈非空，打印每次入栈、弹栈和结算对象。重点观察相等元素、空栈、哨兵和最后一次结算。

\textbf{复杂度变化。} 反复删除最坏 O(\ensuremath{n^{2}})，栈扫描 O(n) 时间，O(n) 空间。

\textbf{迁移追问。} 如果题目改成“若加入通配符星号，需要额外维护可行区间或双栈”，先判断原状态是否仍然覆盖新答案集合，再决定是微调边界、改 value 语义，还是换算法。

\noindent\textbf{LeetCode 32 最长有效括号。}

\textbf{题目说明。} LeetCode 32 最长有效括号 的题面入口要先还原成具体任务：有效括号子串需要知道每段非法边界之后的最早可连接位置。

\textbf{样例入口。} 拿字符串 )()()) 手算。第一个 ) 不是可匹配对象，而是非法边界；后面每次匹配成功，当前有效长度由当前位置减去栈顶边界得出。若栈空，要把当前右括号作为新非法边界压入。

\textbf{暴力基线。} 慢版本让每个位置向左或向右线性寻找边界，或者让每个结构反复等待匹配。它能算出答案，但很多尚未完成的候选会被未来同一个事件一起解决。放到本题就是：先按“有效括号子串需要知道每段非法边界之后的最早可连接位置”完整试慢版本；样例里已经能看到“规律是栈里既保存未匹配左括号下标，也保存最近非法边界”，所以优化前必须先能说清慢版本枚举了哪些候选。

\textbf{暴力过程。} 拿字符串 )()()) 手算。第一个 ) 不是可匹配对象，而是非法边界；后面每次匹配成功，当前有效长度由当前位置减去栈顶边界得出。若栈空，要把当前右括号作为新非法边界压入。

\textbf{重复工作。} 题面模型：有效括号子串需要知道每段非法边界之后的最早可连接位置。暴力版的慢点要落到本题：慢版本会围绕“有效括号子串需要知道每段非法边界之后的最早可连接位置”反复寻找最近边界或最近未完成对象。样例规律是“规律是栈里既保存未匹配左括号下标，也保存最近非法边界”，说明旧候选一旦遇到当前输入就能被批量结算；继续为每个位置单独向左或向右扫描就是浪费。后面的优化只消掉这类重复工作，不能改变答案集合。

\textbf{优化条件。} 本题的关键观察是：栈保存下标，先放入哨兵负一；匹配后用当前下标减栈顶得到长度。这句话必须能翻译成可维护状态；如果题目条件改变后观察不再成立，就不能继续套原来的写法。

\textbf{相邻状态。} 规律是栈里既保存未匹配左括号下标，也保存最近非法边界。把这个特例继续拆开看：栈里的对象都是答案尚未确定的候选；一旦当前元素触发弹栈，被弹出的候选必须已经同时拿到了左边界和右边界，未来元素不再有资格改变它的答案。

\textbf{状态复用。} 把反复寻找最近边界改成维护未完成候选；新元素只和栈顶比较，连续弹出一批已经被当前元素解决的对象。本题落点是：规律是栈里既保存未匹配左括号下标，也保存最近非法边界。手算时只改被新输入影响的那一小部分，而不是回头重算完整候选。

\textbf{指针和字段。} 状态字段要能说清：栈内对象的业务含义、当前输入、弹出时结算的对象、左边界和右边界；栈中存下标还是值要明确。手算时按这个协议跟踪：拿字符串 )()()) 手算。第一个 ) 不是可匹配对象，而是非法边界；后面每次匹配成功，当前有效长度由当前位置减去栈顶边界得出。若栈空，要把当前右括号作为新非法边界压入。然后按协议复查：手算时记录栈内元素的业务含义，不只记录数值。入栈表示答案未定，弹栈表示当前事件已经让它的答案定下来。

\textbf{代码落点。} 弹栈前确认栈非空，弹出后再读取新的左边界；相等元素和哨兵高度要有统一策略。关键代码行必须能逐句翻译成“旧状态怎样变成新状态”，否则就只是模板。

\textbf{正确性。} 证明从弹出时机开始：被弹出的对象在当前时刻答案已经确定，未来元素无法再改变它。

\textbf{边界实验。} 先用空串、全左括号、全右括号、多个断开的合法段做反例检查。实验时准备空串、全左括号、全右括号、多个断开的合法段，打印每次入栈、弹栈和结算对象。重点观察相等元素、空栈、哨兵和最后一次结算。

\textbf{复杂度变化。} 暴力为每个位置寻找边界常见为平方级；单调栈让每个元素最多入栈一次、出栈一次，因此主体扫描为线性时间。

\textbf{迁移追问。} 如果题目改成“DP 写法也可做，但栈写法更直接表达最近非法边界”，先判断原状态是否仍然覆盖新答案集合，再决定是微调边界、改 value 语义，还是换算法。

\noindent\textbf{LeetCode 42 接雨水。}

\textbf{题目说明。} LeetCode 42 接雨水 的题面入口要先还原成具体任务：每个位置的水由左侧最高和右侧最高的较小者决定。

\textbf{样例入口。} 拿高度 [0, 2, 0, 3] 手算，中间的 0 左侧最高是 2，右侧最高是 3，所以能接 2 格水。若用单调栈，遇到 3 时弹出 0，左边界是 2，右边界是 3，宽度为 1，高度差为 2。

\textbf{暴力基线。} 暴力对每个位置 i，向左扫描最高柱 leftMax，向右扫描最高柱 rightMax，当前位置接水为 min(leftMax, rightMax)-height[i]。

\textbf{暴力过程。} 以 [0, 1, 0, 2] 为例，下标 2 左侧最高为 1，右侧最高为 2，所以这里能接 1 格水。

\textbf{重复工作。} 题面模型：每个位置的水由左侧最高和右侧最高的较小者决定。暴力版的慢点要落到本题：浪费在每个位置都重新向两边找最高柱。左右最高值可以前缀后缀预处理，也可以用双指针增量维护。后面的优化只消掉这类重复工作，不能改变答案集合。

\textbf{优化条件。} 本题的关键观察是：前后缀、双指针、单调栈都是不同状态维护方式；要能说明各自结算对象。这句话必须能翻译成可维护状态；如果题目条件改变后观察不再成立，就不能继续套原来的写法。

\textbf{相邻状态。} 规律是每个凹槽必须同时等到左右边界，弹栈那一刻才真正结算水量。把这个特例继续拆开看：栈里的对象都是答案尚未确定的候选；一旦当前元素触发弹栈，被弹出的候选必须已经同时拿到了左边界和右边界，未来元素不再有资格改变它的答案。

\textbf{状态复用。} 双指针维护 leftMax 和 rightMax。每次处理较低一侧，因为该侧的水量已经由它自己的最高值和另一侧不低的边界确定。手算时只改被新输入影响的那一小部分，而不是回头重算完整候选。

\textbf{指针和字段。} 状态字段要能说清：left/right 是未处理区两端，leftMax/rightMax 是两侧已见最高柱，water 累加已确定位置的贡献。手算时按这个协议跟踪：若 height[left] < height[right]，右侧至少存在 height[right] 这个边界，left 位置能接多少只取决于 leftMax。

\textbf{代码落点。} 关键是先更新 leftMax/rightMax，再累加 max(0, leftMax-height[left]) 或 max(0, rightMax-height[right])。移动的是较低一侧。关键代码行必须能逐句翻译成“旧状态怎样变成新状态”，否则就只是模板。

\textbf{正确性。} 较低一侧的对侧边界已经足够高，不会成为短板；未来内部柱子只能影响未处理位置，不会改变当前已结算位置的水量。

\textbf{边界实验。} 先用单调递增、单调递减、平台高度、多个凹槽、数组长度不足三做反例检查。实验时准备单调递增、单调递减、平台高度、多个凹槽、数组长度不足三，打印每次入栈、弹栈和结算对象。重点观察相等元素、空栈、哨兵和最后一次结算。

\textbf{复杂度变化。} 逐点双向扫描 O(\ensuremath{n^{2}})，前后缀 O(n) 空间，双指针 O(n) 时间、O(1) 空间。

\textbf{迁移追问。} 如果题目改成“若柱子变成二维网格，需要最小堆从边界向内扩展”，先判断原状态是否仍然覆盖新答案集合，再决定是微调边界、改 value 语义，还是换算法。

\noindent\textbf{LeetCode 71 简化路径。}

\textbf{题目说明。} LeetCode 71 简化路径 的题面入口要先还原成具体任务：Unix 绝对路径由斜杠分隔目录段，点号段和返回上级目录段需要被规范化。

\textbf{样例入口。} 拿 path=/home//foo/ 手算，连续斜杠产生空段要忽略，home 和 foo 依次入栈，结果是 /home/foo；拿 /../ 手算，根目录上遇到 .. 不能再向上弹。

\textbf{暴力基线。} 慢版本让每个位置向左或向右线性寻找边界，或者让每个结构反复等待匹配。它能算出答案，但很多尚未完成的候选会被未来同一个事件一起解决。放到本题就是：先按“Unix 绝对路径由斜杠分隔目录段，点号段和返回上级目录段需要被规范化”完整试慢版本；样例里已经能看到“规律是按斜杠切分段，普通目录入栈，.. 弹出有效目录，空段和 . 忽略”，所以优化前必须先能说清慢版本枚举了哪些候选。

\textbf{暴力过程。} 拿 path=/home//foo/ 手算，连续斜杠产生空段要忽略，home 和 foo 依次入栈，结果是 /home/foo；拿 /../ 手算，根目录上遇到 .. 不能再向上弹。

\textbf{重复工作。} 题面模型：Unix 绝对路径由斜杠分隔目录段，点号段和返回上级目录段需要被规范化。暴力版的慢点要落到本题：慢版本会围绕“Unix 绝对路径由斜杠分隔目录段，点号段和返回上级目录段需要被规范化”反复寻找最近边界或最近未完成对象。样例规律是“规律是按斜杠切分段，普通目录入栈，.. 弹出有效目录，空段和 . 忽略”，说明旧候选一旦遇到当前输入就能被批量结算；继续为每个位置单独向左或向右扫描就是浪费。后面的优化只消掉这类重复工作，不能改变答案集合。

\textbf{优化条件。} 本题的关键观察是：按斜杠切分路径，普通目录入栈，遇到 .. 弹出一个有效目录，空段和 . 忽略。这句话必须能翻译成可维护状态；如果题目条件改变后观察不再成立，就不能继续套原来的写法。

\textbf{相邻状态。} 规律是按斜杠切分段，普通目录入栈，.. 弹出有效目录，空段和 . 忽略。把这个特例继续拆开看：栈里的对象都是答案尚未确定的候选；一旦当前元素触发弹栈，被弹出的候选必须已经同时拿到了左边界和右边界，未来元素不再有资格改变它的答案。

\textbf{状态复用。} 把反复寻找最近边界改成维护未完成候选；新元素只和栈顶比较，连续弹出一批已经被当前元素解决的对象。本题落点是：规律是按斜杠切分段，普通目录入栈，.. 弹出有效目录，空段和 . 忽略。手算时只改被新输入影响的那一小部分，而不是回头重算完整候选。

\textbf{指针和字段。} 状态字段要能说清：栈内对象的业务含义、当前输入、弹出时结算的对象、左边界和右边界；栈中存下标还是值要明确。手算时按这个协议跟踪：拿 path=/home//foo/ 手算，连续斜杠产生空段要忽略，home 和 foo 依次入栈，结果是 /home/foo；拿 /../ 手算，根目录上遇到 .. 不能再向上弹。然后按协议复查：手算时记录栈内元素的业务含义，不只记录数值。入栈表示答案未定，弹栈表示当前事件已经让它的答案定下来。

\textbf{代码落点。} 弹栈前确认栈非空，弹出后再读取新的左边界；相等元素和哨兵高度要有统一策略。关键代码行必须能逐句翻译成“旧状态怎样变成新状态”，否则就只是模板。

\textbf{正确性。} 证明从弹出时机开始：被弹出的对象在当前时刻答案已经确定，未来元素无法再改变它。

\textbf{边界实验。} 先用连续斜杠、根目录上继续 ..、路径末尾斜杠、目录名含点但不是点号做反例检查。实验时准备连续斜杠、根目录上继续 ..、路径末尾斜杠、目录名含点但不是点号，打印每次入栈、弹栈和结算对象。重点观察相等元素、空栈、哨兵和最后一次结算。

\textbf{复杂度变化。} 暴力为每个位置寻找边界常见为平方级；单调栈让每个元素最多入栈一次、出栈一次，因此主体扫描为线性时间。

\textbf{迁移追问。} 如果题目改成“若路径是相对路径，开头和越过根目录的处理规则要重新定义”，先判断原状态是否仍然覆盖新答案集合，再决定是微调边界、改 value 语义，还是换算法。

\noindent\textbf{LeetCode 84 柱状图中最大的矩形。}

\textbf{题目说明。} LeetCode 84 柱状图中最大的矩形 的题面入口要先还原成具体任务：每根柱子作为最矮高度时，需要知道左右第一个更矮位置。

\textbf{样例入口。} 拿高度 [2, 1, 2] 手算。第一个 2 遇到更矮的 1 时，右边界已经确定在当前位置，左边界是弹出后新栈顶的后一位，所以面积可结算。若一直递增，最后需要哨兵触发结算。

\textbf{暴力基线。} 暴力枚举每根柱子作为矩形最矮高度，再向左右扩展直到遇到更矮柱子，计算最大面积。

\textbf{暴力过程。} 以 heights=[2, 1, 5, 6, 2, 3] 为例，高度 5 的柱子向右可覆盖 6，向左遇到 1 停止，面积可为 5*2=10。

\textbf{重复工作。} 题面模型：每根柱子作为最矮高度时，需要知道左右第一个更矮位置。暴力版的慢点要落到本题：浪费在每根柱子都重新找左右第一个更矮位置。单调递增栈能在遇到更矮柱时一次结算被弹出柱子的右边界。后面的优化只消掉这类重复工作，不能改变答案集合。

\textbf{优化条件。} 本题的关键观察是：单调递增栈在遇到更矮柱时结算被弹出柱子的最大宽度。这句话必须能翻译成可维护状态；如果题目条件改变后观察不再成立，就不能继续套原来的写法。

\textbf{相邻状态。} 规律是栈保存递增高度下标，弹出时才知道该柱子的最大扩展宽度。把这个特例继续拆开看：栈里的对象都是答案尚未确定的候选；一旦当前元素触发弹栈，被弹出的候选必须已经同时拿到了左边界和右边界，未来元素不再有资格改变它的答案。

\textbf{状态复用。} 栈保存下标，且高度单调递增。扫描到更矮高度时，弹出 top，当前下标是右边界，弹出后新的栈顶是左边界。手算时只改被新输入影响的那一小部分，而不是回头重算完整候选。

\textbf{指针和字段。} 状态字段要能说清：top 是被结算柱子，height[top] 是矩形高度，left\_smaller=stack.top，right\_smaller=i，宽度为 i-left\_smaller-1。手算时按这个协议跟踪：扫描到高度 2 时，6 和 5 依次被弹出；弹出 5 时左边界是高度 1 的下标，右边界是当前 2 的下标，面积 10。

\textbf{代码落点。} 关键是两端加哨兵 0，或循环结束后手动清栈。相等高度要有统一策略，常见写法用 >= 弹出保证边界一致。关键代码行必须能逐句翻译成“旧状态怎样变成新状态”，否则就只是模板。

\textbf{正确性。} 柱子被弹出时，左右第一个更矮位置都已经确定，未来更右的柱子不能再扩大以它为最低高度的矩形。

\textbf{边界实验。} 先用递增、递减、相等高度、哨兵零、宽度计算做反例检查。实验时准备递增、递减、相等高度、哨兵零、宽度计算，打印每次入栈、弹栈和结算对象。重点观察相等元素、空栈、哨兵和最后一次结算。

\textbf{复杂度变化。} 暴力 O(\ensuremath{n^{2}})，单调栈每个下标入栈出栈一次，O(n) 时间、O(n) 空间。

\textbf{迁移追问。} 如果题目改成“最大矩形可以作为二维矩阵最大矩形的行压缩子问题”，先判断原状态是否仍然覆盖新答案集合，再决定是微调边界、改 value 语义，还是换算法。

\noindent\textbf{LeetCode 85 最大矩形。}

\textbf{题目说明。} LeetCode 85 最大矩形 的题面入口要先还原成具体任务：每一行都可以把连续的一当成柱状图高度。

\textbf{样例入口。} 拿矩阵每一行向上累计高度手算，某行高度数组若为 [2, 1, 2]，就退化成柱状图最大矩形。

\textbf{暴力基线。} 慢版本让每个位置向左或向右线性寻找边界，或者让每个结构反复等待匹配。它能算出答案，但很多尚未完成的候选会被未来同一个事件一起解决。放到本题就是：先按“每一行都可以把连续的一当成柱状图高度”完整试慢版本；样例里已经能看到“规律是逐行更新每列连续 1 的高度，再对每行高度数组跑单调栈；全零行会把高度清零”，所以优化前必须先能说清慢版本枚举了哪些候选。

\textbf{暴力过程。} 拿矩阵每一行向上累计高度手算，某行高度数组若为 [2, 1, 2]，就退化成柱状图最大矩形。

\textbf{重复工作。} 题面模型：每一行都可以把连续的一当成柱状图高度。暴力版的慢点要落到本题：慢版本会围绕“每一行都可以把连续的一当成柱状图高度”反复寻找最近边界或最近未完成对象。样例规律是“规律是逐行更新每列连续 1 的高度，再对每行高度数组跑单调栈；全零行会把高度清零”，说明旧候选一旦遇到当前输入就能被批量结算；继续为每个位置单独向左或向右扫描就是浪费。后面的优化只消掉这类重复工作，不能改变答案集合。

\textbf{优化条件。} 本题的关键观察是：逐行更新每列高度，然后把当前行转成柱状图最大矩形问题。这句话必须能翻译成可维护状态；如果题目条件改变后观察不再成立，就不能继续套原来的写法。

\textbf{相邻状态。} 规律是逐行更新每列连续 1 的高度，再对每行高度数组跑单调栈；全零行会把高度清零。把这个特例继续拆开看：栈里的对象都是答案尚未确定的候选；一旦当前元素触发弹栈，被弹出的候选必须已经同时拿到了左边界和右边界，未来元素不再有资格改变它的答案。

\textbf{状态复用。} 把反复寻找最近边界改成维护未完成候选；新元素只和栈顶比较，连续弹出一批已经被当前元素解决的对象。本题落点是：规律是逐行更新每列连续 1 的高度，再对每行高度数组跑单调栈；全零行会把高度清零。手算时只改被新输入影响的那一小部分，而不是回头重算完整候选。

\textbf{指针和字段。} 状态字段要能说清：栈内对象的业务含义、当前输入、弹出时结算的对象、左边界和右边界；栈中存下标还是值要明确。手算时按这个协议跟踪：拿矩阵每一行向上累计高度手算，某行高度数组若为 [2, 1, 2]，就退化成柱状图最大矩形。然后按协议复查：手算时记录栈内元素的业务含义，不只记录数值。入栈表示答案未定，弹栈表示当前事件已经让它的答案定下来。

\textbf{代码落点。} 弹栈前确认栈非空，弹出后再读取新的左边界；相等元素和哨兵高度要有统一策略。关键代码行必须能逐句翻译成“旧状态怎样变成新状态”，否则就只是模板。

\textbf{正确性。} 证明从弹出时机开始：被弹出的对象在当前时刻答案已经确定，未来元素无法再改变它。

\textbf{边界实验。} 先用全零、全一、单行、单列、宽度和高度结算做反例检查。实验时准备全零、全一、单行、单列、宽度和高度结算，打印每次入栈、弹栈和结算对象。重点观察相等元素、空栈、哨兵和最后一次结算。

\textbf{复杂度变化。} 暴力为每个位置寻找边界常见为平方级；单调栈让每个元素最多入栈一次、出栈一次，因此主体扫描为线性时间。

\textbf{迁移追问。} 如果题目改成“这是二维问题压成一维单调栈的典型做法”，先判断原状态是否仍然覆盖新答案集合，再决定是微调边界、改 value 语义，还是换算法。

\noindent\textbf{LeetCode 150 逆波兰表达式求值。}

\textbf{题目说明。} LeetCode 150 逆波兰表达式求值 的题面入口要先还原成具体任务：后缀表达式把优先级和括号都编码进 token 顺序。

\textbf{样例入口。} 拿 tokens=[2, 1, +, 3, *] 手算，2 和 1 入栈，遇到 + 弹出右操作数 1 和左操作数 2 得 3，再与 3 相乘得 9。

\textbf{暴力基线。} 慢版本让每个位置向左或向右线性寻找边界，或者让每个结构反复等待匹配。它能算出答案，但很多尚未完成的候选会被未来同一个事件一起解决。放到本题就是：先按“后缀表达式把优先级和括号都编码进 token 顺序”完整试慢版本；样例里已经能看到“规律是后缀表达式遇到运算符时，栈顶两个数已经是它的完整左右操作数，减法和除法顺序不能反”，所以优化前必须先能说清慢版本枚举了哪些候选。

\textbf{暴力过程。} 拿 tokens=[2, 1, +, 3, *] 手算，2 和 1 入栈，遇到 + 弹出右操作数 1 和左操作数 2 得 3，再与 3 相乘得 9。

\textbf{重复工作。} 题面模型：后缀表达式把优先级和括号都编码进 token 顺序。暴力版的慢点要落到本题：慢版本会围绕“后缀表达式把优先级和括号都编码进 token 顺序”反复寻找最近边界或最近未完成对象。样例规律是“规律是后缀表达式遇到运算符时，栈顶两个数已经是它的完整左右操作数，减法和除法顺序不能反”，说明旧候选一旦遇到当前输入就能被批量结算；继续为每个位置单独向左或向右扫描就是浪费。后面的优化只消掉这类重复工作，不能改变答案集合。

\textbf{优化条件。} 本题的关键观察是：遇到数字入栈，遇到运算符弹出右操作数和左操作数计算。这句话必须能翻译成可维护状态；如果题目条件改变后观察不再成立，就不能继续套原来的写法。

\textbf{相邻状态。} 规律是后缀表达式遇到运算符时，栈顶两个数已经是它的完整左右操作数，减法和除法顺序不能反。把这个特例继续拆开看：栈里的对象都是答案尚未确定的候选；一旦当前元素触发弹栈，被弹出的候选必须已经同时拿到了左边界和右边界，未来元素不再有资格改变它的答案。

\textbf{状态复用。} 把反复寻找最近边界改成维护未完成候选；新元素只和栈顶比较，连续弹出一批已经被当前元素解决的对象。本题落点是：规律是后缀表达式遇到运算符时，栈顶两个数已经是它的完整左右操作数，减法和除法顺序不能反。手算时只改被新输入影响的那一小部分，而不是回头重算完整候选。

\textbf{指针和字段。} 状态字段要能说清：栈内对象的业务含义、当前输入、弹出时结算的对象、左边界和右边界；栈中存下标还是值要明确。手算时按这个协议跟踪：拿 tokens=[2, 1, +, 3, *] 手算，2 和 1 入栈，遇到 + 弹出右操作数 1 和左操作数 2 得 3，再与 3 相乘得 9。然后按协议复查：手算时记录栈内元素的业务含义，不只记录数值。入栈表示答案未定，弹栈表示当前事件已经让它的答案定下来。

\textbf{代码落点。} 弹栈前确认栈非空，弹出后再读取新的左边界；相等元素和哨兵高度要有统一策略。关键代码行必须能逐句翻译成“旧状态怎样变成新状态”，否则就只是模板。

\textbf{正确性。} 证明从弹出时机开始：被弹出的对象在当前时刻答案已经确定，未来元素无法再改变它。

\textbf{边界实验。} 先用负数 token、多位数、除法向零截断、减法顺序做反例检查。实验时准备负数 token、多位数、除法向零截断、减法顺序，打印每次入栈、弹栈和结算对象。重点观察相等元素、空栈、哨兵和最后一次结算。

\textbf{复杂度变化。} 暴力为每个位置寻找边界常见为平方级；单调栈让每个元素最多入栈一次、出栈一次，因此主体扫描为线性时间。

\textbf{迁移追问。} 如果题目改成“中缀表达式求值需要另一个运算符栈处理优先级”，先判断原状态是否仍然覆盖新答案集合，再决定是微调边界、改 value 语义，还是换算法。

\noindent\textbf{LeetCode 155 最小栈。}

\textbf{题目说明。} LeetCode 155 最小栈 的题面入口要先还原成具体任务：普通栈只能看到栈顶，不能常数时间知道历史最小值。

\textbf{样例入口。} 拿 push 2、push 0、push 3、push 0 手算，当前最小值序列是 2、0、0、0；弹出最后一个 0 后，最小值仍是 0。

\textbf{暴力基线。} 慢版本让每个位置向左或向右线性寻找边界，或者让每个结构反复等待匹配。它能算出答案，但很多尚未完成的候选会被未来同一个事件一起解决。放到本题就是：先按“普通栈只能看到栈顶，不能常数时间知道历史最小值”完整试慢版本；样例里已经能看到“规律是辅助栈要在每个深度同步保存当前最小值，重复最小值不能只保存一次”，所以优化前必须先能说清慢版本枚举了哪些候选。

\textbf{暴力过程。} 拿 push 2、push 0、push 3、push 0 手算，当前最小值序列是 2、0、0、0；弹出最后一个 0 后，最小值仍是 0。

\textbf{重复工作。} 题面模型：普通栈只能看到栈顶，不能常数时间知道历史最小值。暴力版的慢点要落到本题：慢版本会围绕“普通栈只能看到栈顶，不能常数时间知道历史最小值”反复寻找最近边界或最近未完成对象。样例规律是“规律是辅助栈要在每个深度同步保存当前最小值，重复最小值不能只保存一次”，说明旧候选一旦遇到当前输入就能被批量结算；继续为每个位置单独向左或向右扫描就是浪费。后面的优化只消掉这类重复工作，不能改变答案集合。

\textbf{优化条件。} 本题的关键观察是：辅助栈同步保存每个深度对应的最小值。这句话必须能翻译成可维护状态；如果题目条件改变后观察不再成立，就不能继续套原来的写法。

\textbf{相邻状态。} 规律是辅助栈要在每个深度同步保存当前最小值，重复最小值不能只保存一次。把这个特例继续拆开看：栈里的对象都是答案尚未确定的候选；一旦当前元素触发弹栈，被弹出的候选必须已经同时拿到了左边界和右边界，未来元素不再有资格改变它的答案。

\textbf{状态复用。} 把反复寻找最近边界改成维护未完成候选；新元素只和栈顶比较，连续弹出一批已经被当前元素解决的对象。本题落点是：规律是辅助栈要在每个深度同步保存当前最小值，重复最小值不能只保存一次。手算时只改被新输入影响的那一小部分，而不是回头重算完整候选。

\textbf{指针和字段。} 状态字段要能说清：栈内对象的业务含义、当前输入、弹出时结算的对象、左边界和右边界；栈中存下标还是值要明确。手算时按这个协议跟踪：拿 push 2、push 0、push 3、push 0 手算，当前最小值序列是 2、0、0、0；弹出最后一个 0 后，最小值仍是 0。然后按协议复查：手算时记录栈内元素的业务含义，不只记录数值。入栈表示答案未定，弹栈表示当前事件已经让它的答案定下来。

\textbf{代码落点。} 弹栈前确认栈非空，弹出后再读取新的左边界；相等元素和哨兵高度要有统一策略。关键代码行必须能逐句翻译成“旧状态怎样变成新状态”，否则就只是模板。

\textbf{正确性。} 证明从弹出时机开始：被弹出的对象在当前时刻答案已经确定，未来元素无法再改变它。

\textbf{边界实验。} 先用重复最小值、弹出最小值、空栈操作约定做反例检查。实验时准备重复最小值、弹出最小值、空栈操作约定，打印每次入栈、弹栈和结算对象。重点观察相等元素、空栈、哨兵和最后一次结算。

\textbf{复杂度变化。} 暴力为每个位置寻找边界常见为平方级；单调栈让每个元素最多入栈一次、出栈一次，因此主体扫描为线性时间。

\textbf{迁移追问。} 如果题目改成“也可保存差值压缩空间，但可读性和溢出处理更复杂”，先判断原状态是否仍然覆盖新答案集合，再决定是微调边界、改 value 语义，还是换算法。

\noindent\textbf{LeetCode 224 基本计算器。}

\textbf{题目说明。} LeetCode 224 基本计算器 的题面入口要先还原成具体任务：括号会开启新的表达式上下文，括号结束时要回到上一层。

\textbf{样例入口。} 拿 1-(2+3) 手算，进入括号前当前结果是 1，符号是 -；括号内算出 5 后要乘以上层符号并合并成 -4。

\textbf{暴力基线。} 慢版本让每个位置向左或向右线性寻找边界，或者让每个结构反复等待匹配。它能算出答案，但很多尚未完成的候选会被未来同一个事件一起解决。放到本题就是：先按“括号会开启新的表达式上下文，括号结束时要回到上一层”完整试慢版本；样例里已经能看到“规律是栈保存进入括号前的结果和符号，一元负号和空格要按字符流处理”，所以优化前必须先能说清慢版本枚举了哪些候选。

\textbf{暴力过程。} 拿 1-(2+3) 手算，进入括号前当前结果是 1，符号是 -；括号内算出 5 后要乘以上层符号并合并成 -4。

\textbf{重复工作。} 题面模型：括号会开启新的表达式上下文，括号结束时要回到上一层。暴力版的慢点要落到本题：慢版本会围绕“括号会开启新的表达式上下文，括号结束时要回到上一层”反复寻找最近边界或最近未完成对象。样例规律是“规律是栈保存进入括号前的结果和符号，一元负号和空格要按字符流处理”，说明旧候选一旦遇到当前输入就能被批量结算；继续为每个位置单独向左或向右扫描就是浪费。后面的优化只消掉这类重复工作，不能改变答案集合。

\textbf{优化条件。} 本题的关键观察是：栈保存进入括号前的结果和符号，当前层算完后乘以上层符号并合并。这句话必须能翻译成可维护状态；如果题目条件改变后观察不再成立，就不能继续套原来的写法。

\textbf{相邻状态。} 规律是栈保存进入括号前的结果和符号，一元负号和空格要按字符流处理。把这个特例继续拆开看：栈里的对象都是答案尚未确定的候选；一旦当前元素触发弹栈，被弹出的候选必须已经同时拿到了左边界和右边界，未来元素不再有资格改变它的答案。

\textbf{状态复用。} 把反复寻找最近边界改成维护未完成候选；新元素只和栈顶比较，连续弹出一批已经被当前元素解决的对象。本题落点是：规律是栈保存进入括号前的结果和符号，一元负号和空格要按字符流处理。手算时只改被新输入影响的那一小部分，而不是回头重算完整候选。

\textbf{指针和字段。} 状态字段要能说清：栈内对象的业务含义、当前输入、弹出时结算的对象、左边界和右边界；栈中存下标还是值要明确。手算时按这个协议跟踪：拿 1-(2+3) 手算，进入括号前当前结果是 1，符号是 -；括号内算出 5 后要乘以上层符号并合并成 -4。然后按协议复查：手算时记录栈内元素的业务含义，不只记录数值。入栈表示答案未定，弹栈表示当前事件已经让它的答案定下来。

\textbf{代码落点。} 弹栈前确认栈非空，弹出后再读取新的左边界；相等元素和哨兵高度要有统一策略。关键代码行必须能逐句翻译成“旧状态怎样变成新状态”，否则就只是模板。

\textbf{正确性。} 证明从弹出时机开始：被弹出的对象在当前时刻答案已经确定，未来元素无法再改变它。

\textbf{边界实验。} 先用多位数、前导空格、一元负号、嵌套括号做反例检查。实验时准备多位数、前导空格、一元负号、嵌套括号，打印每次入栈、弹栈和结算对象。重点观察相等元素、空栈、哨兵和最后一次结算。

\textbf{复杂度变化。} 暴力为每个位置寻找边界常见为平方级；单调栈让每个元素最多入栈一次、出栈一次，因此主体扫描为线性时间。

\textbf{迁移追问。} 如果题目改成“基本计算器 II 没有括号但有乘除优先级，需要不同的栈语义”，先判断原状态是否仍然覆盖新答案集合，再决定是微调边界、改 value 语义，还是换算法。

\noindent\textbf{LeetCode 227 基本计算器 II。}

\textbf{题目说明。} LeetCode 227 基本计算器 II 的题面入口要先还原成具体任务：乘除优先级高于加减，可以在读到下一个运算符时结算上一项。

\textbf{样例入口。} 拿 3+2*2 手算，加号前的 3 可以先入栈，乘号遇到 2 时要立刻把栈顶 2 乘以下一个数 2，而不能等到最后按顺序相加。

\textbf{暴力基线。} 慢版本让每个位置向左或向右线性寻找边界，或者让每个结构反复等待匹配。它能算出答案，但很多尚未完成的候选会被未来同一个事件一起解决。放到本题就是：先按“乘除优先级高于加减，可以在读到下一个运算符时结算上一项”完整试慢版本；样例里已经能看到“规律是栈保存已确定符号的项，乘除立即修改栈顶，加减压入正负项”，所以优化前必须先能说清慢版本枚举了哪些候选。

\textbf{暴力过程。} 拿 3+2*2 手算，加号前的 3 可以先入栈，乘号遇到 2 时要立刻把栈顶 2 乘以下一个数 2，而不能等到最后按顺序相加。

\textbf{重复工作。} 题面模型：乘除优先级高于加减，可以在读到下一个运算符时结算上一项。暴力版的慢点要落到本题：慢版本会围绕“乘除优先级高于加减，可以在读到下一个运算符时结算上一项”反复寻找最近边界或最近未完成对象。样例规律是“规律是栈保存已确定符号的项，乘除立即修改栈顶，加减压入正负项”，说明旧候选一旦遇到当前输入就能被批量结算；继续为每个位置单独向左或向右扫描就是浪费。后面的优化只消掉这类重复工作，不能改变答案集合。

\textbf{优化条件。} 本题的关键观察是：栈保存已经确定符号的项，乘除直接修改栈顶，加减压入正负项。这句话必须能翻译成可维护状态；如果题目条件改变后观察不再成立，就不能继续套原来的写法。

\textbf{相邻状态。} 规律是栈保存已确定符号的项，乘除立即修改栈顶，加减压入正负项。把这个特例继续拆开看：栈里的对象都是答案尚未确定的候选；一旦当前元素触发弹栈，被弹出的候选必须已经同时拿到了左边界和右边界，未来元素不再有资格改变它的答案。

\textbf{状态复用。} 把反复寻找最近边界改成维护未完成候选；新元素只和栈顶比较，连续弹出一批已经被当前元素解决的对象。本题落点是：规律是栈保存已确定符号的项，乘除立即修改栈顶，加减压入正负项。手算时只改被新输入影响的那一小部分，而不是回头重算完整候选。

\textbf{指针和字段。} 状态字段要能说清：栈内对象的业务含义、当前输入、弹出时结算的对象、左边界和右边界；栈中存下标还是值要明确。手算时按这个协议跟踪：拿 3+2*2 手算，加号前的 3 可以先入栈，乘号遇到 2 时要立刻把栈顶 2 乘以下一个数 2，而不能等到最后按顺序相加。然后按协议复查：手算时记录栈内元素的业务含义，不只记录数值。入栈表示答案未定，弹栈表示当前事件已经让它的答案定下来。

\textbf{代码落点。} 弹栈前确认栈非空，弹出后再读取新的左边界；相等元素和哨兵高度要有统一策略。关键代码行必须能逐句翻译成“旧状态怎样变成新状态”，否则就只是模板。

\textbf{正确性。} 证明从弹出时机开始：被弹出的对象在当前时刻答案已经确定，未来元素无法再改变它。

\textbf{边界实验。} 先用多位数、空格、除法向零截断、最后一个数字需要结算做反例检查。实验时准备多位数、空格、除法向零截断、最后一个数字需要结算，打印每次入栈、弹栈和结算对象。重点观察相等元素、空栈、哨兵和最后一次结算。

\textbf{复杂度变化。} 暴力为每个位置寻找边界常见为平方级；单调栈让每个元素最多入栈一次、出栈一次，因此主体扫描为线性时间。

\textbf{迁移追问。} 如果题目改成“若加入括号，需要再引入上下文栈或递归下降解析”，先判断原状态是否仍然覆盖新答案集合，再决定是微调边界、改 value 语义，还是换算法。

\noindent\textbf{LeetCode 232 用栈实现队列。}

\textbf{题目说明。} LeetCode 232 用栈实现队列 的题面入口要先还原成具体任务：两个栈分别负责输入顺序和输出顺序。

\textbf{样例入口。} 拿 push 1、push 2、pop 手算，如果输出栈为空，把输入栈整体倒过去得到栈顶 1，正好符合队列先进先出。

\textbf{暴力基线。} 慢版本让每个位置向左或向右线性寻找边界，或者让每个结构反复等待匹配。它能算出答案，但很多尚未完成的候选会被未来同一个事件一起解决。放到本题就是：先按“两个栈分别负责输入顺序和输出顺序”完整试慢版本；样例里已经能看到“规律是只有输出栈为空才搬运，单个元素最多进出两个栈各一次，所以均摊常数”，所以优化前必须先能说清慢版本枚举了哪些候选。

\textbf{暴力过程。} 拿 push 1、push 2、pop 手算，如果输出栈为空，把输入栈整体倒过去得到栈顶 1，正好符合队列先进先出。

\textbf{重复工作。} 题面模型：两个栈分别负责输入顺序和输出顺序。暴力版的慢点要落到本题：慢版本会围绕“两个栈分别负责输入顺序和输出顺序”反复寻找最近边界或最近未完成对象。样例规律是“规律是只有输出栈为空才搬运，单个元素最多进出两个栈各一次，所以均摊常数”，说明旧候选一旦遇到当前输入就能被批量结算；继续为每个位置单独向左或向右扫描就是浪费。后面的优化只消掉这类重复工作，不能改变答案集合。

\textbf{优化条件。} 本题的关键观察是：只有输出栈为空时，才把输入栈整体倒过去，保证均摊常数。这句话必须能翻译成可维护状态；如果题目条件改变后观察不再成立，就不能继续套原来的写法。

\textbf{相邻状态。} 规律是只有输出栈为空才搬运，单个元素最多进出两个栈各一次，所以均摊常数。把这个特例继续拆开看：栈里的对象都是答案尚未确定的候选；一旦当前元素触发弹栈，被弹出的候选必须已经同时拿到了左边界和右边界，未来元素不再有资格改变它的答案。

\textbf{状态复用。} 把反复寻找最近边界改成维护未完成候选；新元素只和栈顶比较，连续弹出一批已经被当前元素解决的对象。本题落点是：规律是只有输出栈为空才搬运，单个元素最多进出两个栈各一次，所以均摊常数。手算时只改被新输入影响的那一小部分，而不是回头重算完整候选。

\textbf{指针和字段。} 状态字段要能说清：栈内对象的业务含义、当前输入、弹出时结算的对象、左边界和右边界；栈中存下标还是值要明确。手算时按这个协议跟踪：拿 push 1、push 2、pop 手算，如果输出栈为空，把输入栈整体倒过去得到栈顶 1，正好符合队列先进先出。然后按协议复查：手算时记录栈内元素的业务含义，不只记录数值。入栈表示答案未定，弹栈表示当前事件已经让它的答案定下来。

\textbf{代码落点。} 弹栈前确认栈非空，弹出后再读取新的左边界；相等元素和哨兵高度要有统一策略。关键代码行必须能逐句翻译成“旧状态怎样变成新状态”，否则就只是模板。

\textbf{正确性。} 证明从弹出时机开始：被弹出的对象在当前时刻答案已经确定，未来元素无法再改变它。

\textbf{边界实验。} 先用连续 peek、空队列、交替 push pop做反例检查。实验时准备连续 peek、空队列、交替 push pop，打印每次入栈、弹栈和结算对象。重点观察相等元素、空栈、哨兵和最后一次结算。

\textbf{复杂度变化。} 暴力为每个位置寻找边界常见为平方级；单调栈让每个元素最多入栈一次、出栈一次，因此主体扫描为线性时间。

\textbf{迁移追问。} 如果题目改成“用队列实现栈则要通过轮转保持最近元素在队首”，先判断原状态是否仍然覆盖新答案集合，再决定是微调边界、改 value 语义，还是换算法。

\noindent\textbf{LeetCode 239 滑动窗口最大值。}

\textbf{题目说明。} LeetCode 239 滑动窗口最大值 的题面入口要先还原成具体任务：每个窗口最大值来自仍未过期且没有被新元素支配的候选。

\textbf{样例入口。} 拿 [1, 3, -1]、窗口 3 手算，队尾的 1 在 3 进入后永远不可能成为最大值，因为 3 更新且更靠右。于是 1 可以弹出。

\textbf{暴力基线。} 暴力枚举每个长度为 k 的窗口，扫描窗口内所有元素求最大值。

\textbf{暴力过程。} 以 [1, 3, -1, -3, 5]、k=3 为例，第一个窗口 [1, 3, -1] 最大值为 3；下一个窗口 [3, -1, -3] 仍要重新扫描会重复看 3 和 -1。

\textbf{重复工作。} 题面模型：每个窗口最大值来自仍未过期且没有被新元素支配的候选。暴力版的慢点要落到本题：浪费在相邻窗口重复求最大值。若新来的数更大且位置更靠右，队尾所有不大于它的旧数未来都不可能再当最大值。后面的优化只消掉这类重复工作，不能改变答案集合。

\textbf{优化条件。} 本题的关键观察是：单调队列保存下标，队尾小于等于新值的候选被弹出。这句话必须能翻译成可维护状态；如果题目条件改变后观察不再成立，就不能继续套原来的写法。

\textbf{相邻状态。} 规律是单调队列保存可能成为当前或未来窗口最大值的下标，过期下标从队首清理。把这个特例继续拆开看：栈里的对象都是答案尚未确定的候选；一旦当前元素触发弹栈，被弹出的候选必须已经同时拿到了左边界和右边界，未来元素不再有资格改变它的答案。

\textbf{状态复用。} 单调递减队列保存下标。入队前从队尾弹出 nums[back]<=nums[i] 的下标；队首若过期 i-k 则弹出；窗口形成后队首就是最大值。手算时只改被新输入影响的那一小部分，而不是回头重算完整候选。

\textbf{指针和字段。} 状态字段要能说清：deque.front 是当前最大值下标，i 是右端，i-k+1 是当前窗口左端。手算时按这个协议跟踪：1 入队后，3 进入时弹出 1，因为 3 更大且更晚；窗口形成后队首 3 给出最大值。

\textbf{代码落点。} 关键是队列保存下标而不是值，这样才能判断过期。顺序通常是清队尾、入队、清队首过期、读答案。关键代码行必须能逐句翻译成“旧状态怎样变成新状态”，否则就只是模板。

\textbf{正确性。} 被队尾新元素弹出的旧元素既更小又更早，在未来任何包含旧元素的窗口里也会包含新元素，因此不会成为最大值。

\textbf{边界实验。} 先用k 为一、k 等于数组长度、重复最大值、队首过期做反例检查。实验时准备k 为一、k 等于数组长度、重复最大值、队首过期，打印每次入栈、弹栈和结算对象。重点观察相等元素、空栈、哨兵和最后一次结算。

\textbf{复杂度变化。} 暴力 O(nk)，单调队列每个下标入队出队一次，O(n) 时间、O(k) 空间。

\textbf{迁移追问。} 如果题目改成“受限子序列和使用同样队列维护最近 k 个 DP 最大值”，先判断原状态是否仍然覆盖新答案集合，再决定是微调边界、改 value 语义，还是换算法。

\noindent\textbf{LeetCode 394 字符串解码。}

\textbf{题目说明。} LeetCode 394 字符串解码 的题面入口要先还原成具体任务：嵌套结构需要保存上一层字符串和重复次数。

\textbf{样例入口。} 拿 3[a2[c]] 手算，遇到左括号前保存当前字符串和重复次数；遇到右括号时弹出上一层，把当前 cc 重复 2 次拼成 acc，再被外层重复 3 次。

\textbf{暴力基线。} 慢版本让每个位置向左或向右线性寻找边界，或者让每个结构反复等待匹配。它能算出答案，但很多尚未完成的候选会被未来同一个事件一起解决。放到本题就是：先按“嵌套结构需要保存上一层字符串和重复次数”完整试慢版本；样例里已经能看到“规律是栈保存进入括号前的上下文，右括号触发一次完整展开”，所以优化前必须先能说清慢版本枚举了哪些候选。

\textbf{暴力过程。} 拿 3[a2[c]] 手算，遇到左括号前保存当前字符串和重复次数；遇到右括号时弹出上一层，把当前 cc 重复 2 次拼成 acc，再被外层重复 3 次。

\textbf{重复工作。} 题面模型：嵌套结构需要保存上一层字符串和重复次数。暴力版的慢点要落到本题：慢版本会围绕“嵌套结构需要保存上一层字符串和重复次数”反复寻找最近边界或最近未完成对象。样例规律是“规律是栈保存进入括号前的上下文，右括号触发一次完整展开”，说明旧候选一旦遇到当前输入就能被批量结算；继续为每个位置单独向左或向右扫描就是浪费。后面的优化只消掉这类重复工作，不能改变答案集合。

\textbf{优化条件。} 本题的关键观察是：遇到左括号入栈当前状态，右括号弹出并拼接展开。这句话必须能翻译成可维护状态；如果题目条件改变后观察不再成立，就不能继续套原来的写法。

\textbf{相邻状态。} 规律是栈保存进入括号前的上下文，右括号触发一次完整展开。把这个特例继续拆开看：栈里的对象都是答案尚未确定的候选；一旦当前元素触发弹栈，被弹出的候选必须已经同时拿到了左边界和右边界，未来元素不再有资格改变它的答案。

\textbf{状态复用。} 把反复寻找最近边界改成维护未完成候选；新元素只和栈顶比较，连续弹出一批已经被当前元素解决的对象。本题落点是：规律是栈保存进入括号前的上下文，右括号触发一次完整展开。手算时只改被新输入影响的那一小部分，而不是回头重算完整候选。

\textbf{指针和字段。} 状态字段要能说清：栈内对象的业务含义、当前输入、弹出时结算的对象、左边界和右边界；栈中存下标还是值要明确。手算时按这个协议跟踪：拿 3[a2[c]] 手算，遇到左括号前保存当前字符串和重复次数；遇到右括号时弹出上一层，把当前 cc 重复 2 次拼成 acc，再被外层重复 3 次。然后按协议复查：手算时记录栈内元素的业务含义，不只记录数值。入栈表示答案未定，弹栈表示当前事件已经让它的答案定下来。

\textbf{代码落点。} 弹栈前确认栈非空，弹出后再读取新的左边界；相等元素和哨兵高度要有统一策略。关键代码行必须能逐句翻译成“旧状态怎样变成新状态”，否则就只是模板。

\textbf{正确性。} 证明从弹出时机开始：被弹出的对象在当前时刻答案已经确定，未来元素无法再改变它。

\textbf{边界实验。} 先用多位数字、嵌套、空片段、数字后必须跟括号做反例检查。实验时准备多位数字、嵌套、空片段、数字后必须跟括号，打印每次入栈、弹栈和结算对象。重点观察相等元素、空栈、哨兵和最后一次结算。

\textbf{复杂度变化。} 暴力为每个位置寻找边界常见为平方级；单调栈让每个元素最多入栈一次、出栈一次，因此主体扫描为线性时间。

\textbf{迁移追问。} 如果题目改成“递归下降解析更通用，但迭代栈更符合工程栈控制”，先判断原状态是否仍然覆盖新答案集合，再决定是微调边界、改 value 语义，还是换算法。

\noindent\textbf{LeetCode 402 移掉 K 位数字。}

\textbf{题目说明。} LeetCode 402 移掉 K 位数字 的题面入口要先还原成具体任务：要删除 k 位使剩余数字最小，本质是让高位尽量小。

\textbf{样例入口。} 拿 1432219、k=3 手算，遇到 3 后面的 2 时，前面的 4 和 3 都比 2 大且可删，弹出能让高位更小；最后去掉前导零。

\textbf{暴力基线。} 慢版本让每个位置向左或向右线性寻找边界，或者让每个结构反复等待匹配。它能算出答案，但很多尚未完成的候选会被未来同一个事件一起解决。放到本题就是：先按“要删除 k 位使剩余数字最小，本质是让高位尽量小”完整试慢版本；样例里已经能看到“规律是单调递增栈让更小数字尽量靠前，k 未用完时从尾部继续删”，所以优化前必须先能说清慢版本枚举了哪些候选。

\textbf{暴力过程。} 拿 1432219、k=3 手算，遇到 3 后面的 2 时，前面的 4 和 3 都比 2 大且可删，弹出能让高位更小；最后去掉前导零。

\textbf{重复工作。} 题面模型：要删除 k 位使剩余数字最小，本质是让高位尽量小。暴力版的慢点要落到本题：慢版本会围绕“要删除 k 位使剩余数字最小，本质是让高位尽量小”反复寻找最近边界或最近未完成对象。样例规律是“规律是单调递增栈让更小数字尽量靠前，k 未用完时从尾部继续删”，说明旧候选一旦遇到当前输入就能被批量结算；继续为每个位置单独向左或向右扫描就是浪费。后面的优化只消掉这类重复工作，不能改变答案集合。

\textbf{优化条件。} 本题的关键观察是：单调递增栈中遇到更小数字时，弹出前面较大的数字并消耗删除次数。这句话必须能翻译成可维护状态；如果题目条件改变后观察不再成立，就不能继续套原来的写法。

\textbf{相邻状态。} 规律是单调递增栈让更小数字尽量靠前，k 未用完时从尾部继续删。把这个特例继续拆开看：栈里的对象都是答案尚未确定的候选；一旦当前元素触发弹栈，被弹出的候选必须已经同时拿到了左边界和右边界，未来元素不再有资格改变它的答案。

\textbf{状态复用。} 把反复寻找最近边界改成维护未完成候选；新元素只和栈顶比较，连续弹出一批已经被当前元素解决的对象。本题落点是：规律是单调递增栈让更小数字尽量靠前，k 未用完时从尾部继续删。手算时只改被新输入影响的那一小部分，而不是回头重算完整候选。

\textbf{指针和字段。} 状态字段要能说清：栈内对象的业务含义、当前输入、弹出时结算的对象、左边界和右边界；栈中存下标还是值要明确。手算时按这个协议跟踪：拿 1432219、k=3 手算，遇到 3 后面的 2 时，前面的 4 和 3 都比 2 大且可删，弹出能让高位更小；最后去掉前导零。然后按协议复查：手算时记录栈内元素的业务含义，不只记录数值。入栈表示答案未定，弹栈表示当前事件已经让它的答案定下来。

\textbf{代码落点。} 弹栈前确认栈非空，弹出后再读取新的左边界；相等元素和哨兵高度要有统一策略。关键代码行必须能逐句翻译成“旧状态怎样变成新状态”，否则就只是模板。

\textbf{正确性。} 证明从弹出时机开始：被弹出的对象在当前时刻答案已经确定，未来元素无法再改变它。

\textbf{边界实验。} 先用前导零、k 未用完、k 等于长度、数字已经递增做反例检查。实验时准备前导零、k 未用完、k 等于长度、数字已经递增，打印每次入栈、弹栈和结算对象。重点观察相等元素、空栈、哨兵和最后一次结算。

\textbf{复杂度变化。} 暴力为每个位置寻找边界常见为平方级；单调栈让每个元素最多入栈一次、出栈一次，因此主体扫描为线性时间。

\textbf{迁移追问。} 如果题目改成“去除重复字母也用单调栈，但还要保证每个字符保留一次”，先判断原状态是否仍然覆盖新答案集合，再决定是微调边界、改 value 语义，还是换算法。

\noindent\textbf{LeetCode 496 下一个更大元素 I。}

\textbf{题目说明。} LeetCode 496 下一个更大元素 I 的题面入口要先还原成具体任务：第二个数组提供每个值右侧第一个更大值的全局关系。

\textbf{样例入口。} 拿 nums2=[1, 3, 4, 2] 手算，1 的下一个更大是 3，3 的下一个更大是 4，4 没有。

\textbf{暴力基线。} 慢版本让每个位置向左或向右线性寻找边界，或者让每个结构反复等待匹配。它能算出答案，但很多尚未完成的候选会被未来同一个事件一起解决。放到本题就是：先按“第二个数组提供每个值右侧第一个更大值的全局关系”完整试慢版本；样例里已经能看到“规律是单调栈在扫描 nums2 时建立值到下一个更大值的映射，再回答 nums1 查询”，所以优化前必须先能说清慢版本枚举了哪些候选。

\textbf{暴力过程。} 拿 nums2=[1, 3, 4, 2] 手算，1 的下一个更大是 3，3 的下一个更大是 4，4 没有。

\textbf{重复工作。} 题面模型：第二个数组提供每个值右侧第一个更大值的全局关系。暴力版的慢点要落到本题：慢版本会围绕“第二个数组提供每个值右侧第一个更大值的全局关系”反复寻找最近边界或最近未完成对象。样例规律是“规律是单调栈在扫描 nums2 时建立值到下一个更大值的映射，再回答 nums1 查询”，说明旧候选一旦遇到当前输入就能被批量结算；继续为每个位置单独向左或向右扫描就是浪费。后面的优化只消掉这类重复工作，不能改变答案集合。

\textbf{优化条件。} 本题的关键观察是：单调栈预处理值到下一个更大值映射，再回答第一个数组查询。这句话必须能翻译成可维护状态；如果题目条件改变后观察不再成立，就不能继续套原来的写法。

\textbf{相邻状态。} 规律是单调栈在扫描 nums2 时建立值到下一个更大值的映射，再回答 nums1 查询。把这个特例继续拆开看：栈里的对象都是答案尚未确定的候选；一旦当前元素触发弹栈，被弹出的候选必须已经同时拿到了左边界和右边界，未来元素不再有资格改变它的答案。

\textbf{状态复用。} 把反复寻找最近边界改成维护未完成候选；新元素只和栈顶比较，连续弹出一批已经被当前元素解决的对象。本题落点是：规律是单调栈在扫描 nums2 时建立值到下一个更大值的映射，再回答 nums1 查询。手算时只改被新输入影响的那一小部分，而不是回头重算完整候选。

\textbf{指针和字段。} 状态字段要能说清：栈内对象的业务含义、当前输入、弹出时结算的对象、左边界和右边界；栈中存下标还是值要明确。手算时按这个协议跟踪：拿 nums2=[1, 3, 4, 2] 手算，1 的下一个更大是 3，3 的下一个更大是 4，4 没有。然后按协议复查：手算时记录栈内元素的业务含义，不只记录数值。入栈表示答案未定，弹栈表示当前事件已经让它的答案定下来。

\textbf{代码落点。} 弹栈前确认栈非空，弹出后再读取新的左边界；相等元素和哨兵高度要有统一策略。关键代码行必须能逐句翻译成“旧状态怎样变成新状态”，否则就只是模板。

\textbf{正确性。} 证明从弹出时机开始：被弹出的对象在当前时刻答案已经确定，未来元素无法再改变它。

\textbf{边界实验。} 先用不存在更大值、递减数组、值唯一假设做反例检查。实验时准备不存在更大值、递减数组、值唯一假设，打印每次入栈、弹栈和结算对象。重点观察相等元素、空栈、哨兵和最后一次结算。

\textbf{复杂度变化。} 暴力为每个位置寻找边界常见为平方级；单调栈让每个元素最多入栈一次、出栈一次，因此主体扫描为线性时间。

\textbf{迁移追问。} 如果题目改成“若值不唯一，应按下标而不是值建映射”，先判断原状态是否仍然覆盖新答案集合，再决定是微调边界、改 value 语义，还是换算法。

\noindent\textbf{LeetCode 503 下一个更大元素 II。}

\textbf{题目说明。} LeetCode 503 下一个更大元素 II 的题面入口要先还原成具体任务：环形数组可以通过遍历两倍下标模拟绕回。

\textbf{样例入口。} 拿 [1, 2, 1] 手算，第一个 1 的下一个更大是 2，最后一个 1 环形回到前面也能找到 2。

\textbf{暴力基线。} 慢版本让每个位置向左或向右线性寻找边界，或者让每个结构反复等待匹配。它能算出答案，但很多尚未完成的候选会被未来同一个事件一起解决。放到本题就是：先按“环形数组可以通过遍历两倍下标模拟绕回”完整试慢版本；样例里已经能看到“规律是扫描两遍模拟环形，第一遍入栈，第二遍只帮助未解决下标结算”，所以优化前必须先能说清慢版本枚举了哪些候选。

\textbf{暴力过程。} 拿 [1, 2, 1] 手算，第一个 1 的下一个更大是 2，最后一个 1 环形回到前面也能找到 2。

\textbf{重复工作。} 题面模型：环形数组可以通过遍历两倍下标模拟绕回。暴力版的慢点要落到本题：慢版本会围绕“环形数组可以通过遍历两倍下标模拟绕回”反复寻找最近边界或最近未完成对象。样例规律是“规律是扫描两遍模拟环形，第一遍入栈，第二遍只帮助未解决下标结算”，说明旧候选一旦遇到当前输入就能被批量结算；继续为每个位置单独向左或向右扫描就是浪费。后面的优化只消掉这类重复工作，不能改变答案集合。

\textbf{优化条件。} 本题的关键观察是：第一遍入栈，第二遍只帮助未解决下标找到答案。这句话必须能翻译成可维护状态；如果题目条件改变后观察不再成立，就不能继续套原来的写法。

\textbf{相邻状态。} 规律是扫描两遍模拟环形，第一遍入栈，第二遍只帮助未解决下标结算。把这个特例继续拆开看：栈里的对象都是答案尚未确定的候选；一旦当前元素触发弹栈，被弹出的候选必须已经同时拿到了左边界和右边界，未来元素不再有资格改变它的答案。

\textbf{状态复用。} 把反复寻找最近边界改成维护未完成候选；新元素只和栈顶比较，连续弹出一批已经被当前元素解决的对象。本题落点是：规律是扫描两遍模拟环形，第一遍入栈，第二遍只帮助未解决下标结算。手算时只改被新输入影响的那一小部分，而不是回头重算完整候选。

\textbf{指针和字段。} 状态字段要能说清：栈内对象的业务含义、当前输入、弹出时结算的对象、左边界和右边界；栈中存下标还是值要明确。手算时按这个协议跟踪：拿 [1, 2, 1] 手算，第一个 1 的下一个更大是 2，最后一个 1 环形回到前面也能找到 2。然后按协议复查：手算时记录栈内元素的业务含义，不只记录数值。入栈表示答案未定，弹栈表示当前事件已经让它的答案定下来。

\textbf{代码落点。} 弹栈前确认栈非空，弹出后再读取新的左边界；相等元素和哨兵高度要有统一策略。关键代码行必须能逐句翻译成“旧状态怎样变成新状态”，否则就只是模板。

\textbf{正确性。} 证明从弹出时机开始：被弹出的对象在当前时刻答案已经确定，未来元素无法再改变它。

\textbf{边界实验。} 先用全递减、全相等、长度一、取模下标做反例检查。实验时准备全递减、全相等、长度一、取模下标，打印每次入栈、弹栈和结算对象。重点观察相等元素、空栈、哨兵和最后一次结算。

\textbf{复杂度变化。} 暴力为每个位置寻找边界常见为平方级；单调栈让每个元素最多入栈一次、出栈一次，因此主体扫描为线性时间。

\textbf{迁移追问。} 如果题目改成“不要在第二遍重复入栈，否则可能死循环或重复计算”，先判断原状态是否仍然覆盖新答案集合，再决定是微调边界、改 value 语义，还是换算法。

\noindent\textbf{LeetCode 739 每日温度。}

\textbf{题目说明。} LeetCode 739 每日温度 的题面入口要先还原成具体任务：每一天等待右侧第一个更高温度。

\textbf{样例入口。} 拿 [73, 74, 75, 71, 69, 72, 76, 73] 手算，73 遇到 74 时等待 1 天；71 要等到 72 才解决。

\textbf{暴力基线。} 慢版本让每个位置向左或向右线性寻找边界，或者让每个结构反复等待匹配。它能算出答案，但很多尚未完成的候选会被未来同一个事件一起解决。放到本题就是：先按“每一天等待右侧第一个更高温度”完整试慢版本；样例里已经能看到“规律是单调递减栈保存尚未等到更高温的下标，遇到更高温时连续弹栈结算天数”，所以优化前必须先能说清慢版本枚举了哪些候选。

\textbf{暴力过程。} 拿 [73, 74, 75, 71, 69, 72, 76, 73] 手算，73 遇到 74 时等待 1 天；71 要等到 72 才解决。

\textbf{重复工作。} 题面模型：每一天等待右侧第一个更高温度。暴力版的慢点要落到本题：慢版本会围绕“每一天等待右侧第一个更高温度”反复寻找最近边界或最近未完成对象。样例规律是“规律是单调递减栈保存尚未等到更高温的下标，遇到更高温时连续弹栈结算天数”，说明旧候选一旦遇到当前输入就能被批量结算；继续为每个位置单独向左或向右扫描就是浪费。后面的优化只消掉这类重复工作，不能改变答案集合。

\textbf{优化条件。} 本题的关键观察是：单调递减栈保存尚未找到答案的下标。这句话必须能翻译成可维护状态；如果题目条件改变后观察不再成立，就不能继续套原来的写法。

\textbf{相邻状态。} 规律是单调递减栈保存尚未等到更高温的下标，遇到更高温时连续弹栈结算天数。把这个特例继续拆开看：栈里的对象都是答案尚未确定的候选；一旦当前元素触发弹栈，被弹出的候选必须已经同时拿到了左边界和右边界，未来元素不再有资格改变它的答案。

\textbf{状态复用。} 把反复寻找最近边界改成维护未完成候选；新元素只和栈顶比较，连续弹出一批已经被当前元素解决的对象。本题落点是：规律是单调递减栈保存尚未等到更高温的下标，遇到更高温时连续弹栈结算天数。手算时只改被新输入影响的那一小部分，而不是回头重算完整候选。

\textbf{指针和字段。} 状态字段要能说清：栈内对象的业务含义、当前输入、弹出时结算的对象、左边界和右边界；栈中存下标还是值要明确。手算时按这个协议跟踪：拿 [73, 74, 75, 71, 69, 72, 76, 73] 手算，73 遇到 74 时等待 1 天；71 要等到 72 才解决。然后按协议复查：手算时记录栈内元素的业务含义，不只记录数值。入栈表示答案未定，弹栈表示当前事件已经让它的答案定下来。

\textbf{代码落点。} 弹栈前确认栈非空，弹出后再读取新的左边界；相等元素和哨兵高度要有统一策略。关键代码行必须能逐句翻译成“旧状态怎样变成新状态”，否则就只是模板。

\textbf{正确性。} 证明从弹出时机开始：被弹出的对象在当前时刻答案已经确定，未来元素无法再改变它。

\textbf{边界实验。} 先用没有更高温、相等温度、递增、递减做反例检查。实验时准备没有更高温、相等温度、递增、递减，打印每次入栈、弹栈和结算对象。重点观察相等元素、空栈、哨兵和最后一次结算。

\textbf{复杂度变化。} 暴力为每个位置寻找边界常见为平方级；单调栈让每个元素最多入栈一次、出栈一次，因此主体扫描为线性时间。

\textbf{迁移追问。} 如果题目改成“下一个更大元素与它同型，只是输出值或距离不同”，先判断原状态是否仍然覆盖新答案集合，再决定是微调边界、改 value 语义，还是换算法。

\noindent\textbf{LeetCode 895 最大频率栈。}

\textbf{题目说明。} LeetCode 895 最大频率栈 的题面入口要先还原成具体任务：弹出时既要频率最高，又要在同频中最近入栈。

\textbf{样例入口。} 拿 push 5, 7, 5, 7, 4, 5 后 pop 手算，频率最高是 5，先弹 5；下一次 5 和 7 同频，弹最近达到该频率的 7。

\textbf{暴力基线。} 慢版本让每个位置向左或向右线性寻找边界，或者让每个结构反复等待匹配。它能算出答案，但很多尚未完成的候选会被未来同一个事件一起解决。放到本题就是：先按“弹出时既要频率最高，又要在同频中最近入栈”完整试慢版本；样例里已经能看到“规律是值到频率、频率到栈、当前最大频率三者共同维护”，所以优化前必须先能说清慢版本枚举了哪些候选。

\textbf{暴力过程。} 拿 push 5, 7, 5, 7, 4, 5 后 pop 手算，频率最高是 5，先弹 5；下一次 5 和 7 同频，弹最近达到该频率的 7。

\textbf{重复工作。} 题面模型：弹出时既要频率最高，又要在同频中最近入栈。暴力版的慢点要落到本题：慢版本会围绕“弹出时既要频率最高，又要在同频中最近入栈”反复寻找最近边界或最近未完成对象。样例规律是“规律是值到频率、频率到栈、当前最大频率三者共同维护”，说明旧候选一旦遇到当前输入就能被批量结算；继续为每个位置单独向左或向右扫描就是浪费。后面的优化只消掉这类重复工作，不能改变答案集合。

\textbf{优化条件。} 本题的关键观察是：哈希表记录值频率，频率到栈记录达到该频率的入栈顺序，维护当前最大频率。这句话必须能翻译成可维护状态；如果题目条件改变后观察不再成立，就不能继续套原来的写法。

\textbf{相邻状态。} 规律是值到频率、频率到栈、当前最大频率三者共同维护。把这个特例继续拆开看：栈里的对象都是答案尚未确定的候选；一旦当前元素触发弹栈，被弹出的候选必须已经同时拿到了左边界和右边界，未来元素不再有资格改变它的答案。

\textbf{状态复用。} 把反复寻找最近边界改成维护未完成候选；新元素只和栈顶比较，连续弹出一批已经被当前元素解决的对象。本题落点是：规律是值到频率、频率到栈、当前最大频率三者共同维护。手算时只改被新输入影响的那一小部分，而不是回头重算完整候选。

\textbf{指针和字段。} 状态字段要能说清：栈内对象的业务含义、当前输入、弹出时结算的对象、左边界和右边界；栈中存下标还是值要明确。手算时按这个协议跟踪：拿 push 5, 7, 5, 7, 4, 5 后 pop 手算，频率最高是 5，先弹 5；下一次 5 和 7 同频，弹最近达到该频率的 7。然后按协议复查：手算时记录栈内元素的业务含义，不只记录数值。入栈表示答案未定，弹栈表示当前事件已经让它的答案定下来。

\textbf{代码落点。} 弹栈前确认栈非空，弹出后再读取新的左边界；相等元素和哨兵高度要有统一策略。关键代码行必须能逐句翻译成“旧状态怎样变成新状态”，否则就只是模板。

\textbf{正确性。} 证明从弹出时机开始：被弹出的对象在当前时刻答案已经确定，未来元素无法再改变它。

\textbf{边界实验。} 先用重复 push、pop 后最大频率下降、同频最近性、空栈不在题意内做反例检查。实验时准备重复 push、pop 后最大频率下降、同频最近性、空栈不在题意内，打印每次入栈、弹栈和结算对象。重点观察相等元素、空栈、哨兵和最后一次结算。

\textbf{复杂度变化。} 暴力为每个位置寻找边界常见为平方级；单调栈让每个元素最多入栈一次、出栈一次，因此主体扫描为线性时间。

\textbf{迁移追问。} 如果题目改成“它是栈和哈希表按频率维度组合的设计题”，先判断原状态是否仍然覆盖新答案集合，再决定是微调边界、改 value 语义，还是换算法。

\subsection{实验复盘}

追问通常会问栈里为什么存下标、相等元素如何处理、弹出时到底结算了谁。请准备用一个严格递增和一个严格递减样例解释入栈出栈次数。

复盘本节时先从 LeetCode 20 有效的括号、LeetCode 32 最长有效括号、LeetCode 42 接雨水 开始：每道题都先手写一个能覆盖全部候选的暴力版，再指出优化结构具体保存了哪一段历史信息，最后用相邻案例比较为什么不能互换解法。
